% !PRINTED 384
\weberpage{384}
%% [running head, not body text -- commented out 2026-08-07] Page furniture (running
%% head, folio, signature) is excluded from this edition by standing policy.
%% 384 \hfill Fünfzehnter Abschnitt. \hfill § 104.

%<<CONT
quadratischen Teiler von $m$, den man wieder in zwei teilerfremde
quadratische Faktoren zerlegt, von denen der eine zu $n$, der andere
zu $n'$ genommen wird, und dadurch ändern sich $(\delta, n)$ und
$(\delta', n')$ nicht.

Man kann hiernach einen Charakter $\chi(\delta, n)$ einer beliebigen
zu $Q$ teilerfremden Zahl folgendermaßen bestimmen.

Man setze:

\weq{18}{n = \varepsilon n'}

und verstehe unter $\varepsilon$ das Produkt der in $n$ aufgehenden Potenzen
solcher Primzahlen, die zugleich in $\delta$ aufgehen; dann ist $\varepsilon$ relativ
prim zu $\delta'$ und $n'$ relativ prim zu $\delta$, und man setze:

\weq{19}{\chi(\delta, n) = (\delta', \varepsilon)(\delta, n'),}

wodurch $\chi(\delta, n)$ eindeutig und immer von Null verschieden bestimmt ist. Ist dann $(a, b, c)$ ein Repräsentant der Klasse $A$,
so ist

\weq{20}{\chi(\delta, ax^2 + bxy + cy^2) = \chi(\delta, A),}

worin $x, y$ beliebige ganze Zahlen mit oder ohne gemeinsamen
Teiler sind, für die $ax^2 + bxy + cy^2$ teilerfremd zu $Q$ wird.

Über das so definierte Symbol $\chi(\delta, n)$ gilt eine Reihe von
Sätzen, die wir noch ableiten müssen.

6. \gesperrt{Sind $n_1$ und $n_2$ irgend zwei zu $Q$ teilerfremde
Zahlen, so ist}

\weq{21}{\chi(\delta, n_1)\chi(\delta, n_2) = \chi(\delta, n_1n_2).}

Denn setzen wir nach (18)
\[
n_1 = \varepsilon_1 n_1', \quad n_2 = \varepsilon_2 n_2',
\]
so ist
\[
n_1 n_2 = \varepsilon_1 \varepsilon_2 n_1' n_2',
\]
und dies geht in (18) über, wenn man $n = n_1n_2$, $\varepsilon = \varepsilon_1\varepsilon_2$,
$n' = n_1'n_2'$ setzt. Also folgt (21) aus § 103 (8).

7. \gesperrt{Läßt man $e$ die Reihe der Divisoren von $n$ durchlaufen, und setzt $n = ee'$, so ist}

\weq{22}{\chi(\delta, n)\overset{e}{\Sigma}(D, e) = \overset{e}{\Sigma}(\delta, e)(\delta', e').}

Wenn die Formel (22) für $n = n_1$ und $n = n_2$ gilt, so gilt
sie, unter der Voraussetzung, daß $n_1, n_2$ relativ prim sind, auch
für $n = n_1n_2$ [nach (21)]; sie braucht also nur noch für eine
Primzahlpotenz $n = p^k$ bewiesen zu werden, wenn $p$ nicht in $Q$
%CONT>>
